% Generated draft: requires mathematical and source review.
\documentclass[11pt]{article}
\usepackage[margin=25mm]{geometry}
\usepackage{fontspec}
\setmainfont{Latin Modern Roman}
\usepackage{amsmath,amssymb,mathtools,unicode-math}
\setmathfont{Latin Modern Math}
\usepackage{enumitem,needspace,xurl,hyperref}
\setlength{\parindent}{0pt}
\setlength{\parskip}{5pt}
\setlength{\emergencystretch}{3em}
\newcommand{\N}{\mathbb{N}}
\newcommand{\Z}{\mathbb{Z}}
\newcommand{\Q}{\mathbb{Q}}
\newcommand{\R}{\mathbb{R}}
\newcommand{\C}{\mathbb{C}}
\newcommand{\F}{\mathbb{F}}
\newcommand{\A}{\mathbb{A}}
\newcommand{\PP}{\mathbb{P}}
\newcommand{\E}{\mathbb{E}}
\newcommand{\eps}{\varepsilon}
\newcommand{\dd}{\,\mathrm{d}}
\begin{document}
\section*{UnsolvedMath ID: 30004561}
\textbf{Source problem code:} catalogue entry
\subsection{Definitions and mathematical statement}
A dimer configuration on a graph \(G=(V,E)\) is a perfect matching: a subset \(M\subseteq E\) such that each vertex is incident to exactly one edge of \(M\). The source invokes classical Kenyon--Okounkov limit-shape theory but does not specify a dimer model, its weights, its height-function convention, or its boundary data.

For terminology, let \((\mathcal L,g)\) be a Lorentzian manifold and let \(X:\Sigma\to\mathcal L\) be a smooth immersion of a two-dimensional surface with nondegenerate induced metric \(h=X^*g\). Writing \(\mathrm{II}\) for its second fundamental form, vanishing mean curvature means
\[
H_X:=\frac12\operatorname{tr}_{h}\mathrm{II}=0.
\]
This gives the usual differential-geometric meaning of minimality where the induced metric is nondegenerate. The supplied record does not specify the ambient Lorentzian space, the causal type of the surface, or the treatment of singular or degenerate points; this definition is therefore terminology, not an additional hypothesis attributed to the source.

\textbf{Source research question.} What is the correspondence between classical Kenyon--Okounkov theory and Lorentz-minimal surfaces, and how does local graph structure lead to conjectural Lorentz-minimality of scaling-limit surfaces?

More explicitly, the requested target is to identify the correspondence and the local graph hypotheses under which the associated scaling-limit surfaces should have vanishing mean curvature in the intended Lorentzian geometry. The source does not define the surface construction, the admissible graph sequences, the scaling and convergence, the boundary conditions, or the meaning of correspondence. Consequently, it does not determine a fully quantified conjecture; these unresolved choices must be specified before a precise general assertion can be formulated.
\subsection{Short English statement}
The problem asks how Kenyon--Okounkov dimer limit-shape theory corresponds to Lorentzian surfaces with zero mean curvature, and which local graph properties explain the conjectured minimality of scaling limits. The supplied question leaves the models, geometric construction, convergence, and boundary assumptions unspecified, so this draft records a research target requiring further formulation.
\subsection{Sources}
\begin{itemize}
\item Local UnsolvedMath registry, \texttt{data/unsolvedmath/problems.json}, permanent ID \(30004561\), source code \texttt{catalogue entry}, title ``Dimer Theory and Lorentz-Minimal Scaling-Limit Surfaces.'' The supplied record belongs to the Oberwolfach Reports Open Problems collection. No individual proposer is identified.
\item Supplied report attribution: \emph{Mini-Workshop: Dimers, Ising and Spanning Trees beyond the Critical Isoradial Case} (2021), \emph{Oberwolfach Reports}, pp.~1759--1789, DOI \texttt{10.4171/owr/2020/35}. \url{https://doi.org/10.4171/owr/2020/35}. The registry records proposed year 2020.
\item Supplied literature source: Kenyon and Okounkov, ``limit shapes and complex Burgers.'' \url{https://arxiv.org/abs/math-ph/0507007}.
\item Supplied literature source: Chelkak, Laslier and Russkikh, ``s-embeddings.'' \url{https://arxiv.org/abs/2109.06272}.
\item Supplied literature source: ``OWR 2020/35 report.'' \url{https://doi.org/10.4171/owr/2020/35}.
\item Literature assessment dated \texttt{2026-08-22}: the supplied curation labels the problem open and states that classical Kenyon--Okounkov limit-shape theory and Lorentz-minimal-surface formulations are related in special dimer and embedding settings, while no general correspondence answering the question has been established. This is supplied dated curation derived from the OWR record, not an independently verified literature assessment. Its time-sensitive status and cited claims require independent verification.
\end{itemize}
\end{document}
